\documentclass{mcmthesis}
\usepackage{hyperref} % 引入 hyperref 宏包，用于添加超链接功能
\usepackage{tocloft}
\mcmsetup{tcn = 12345, problem = C, titlepage = true}
\mcmsetup{sheet = true}

\usepackage{blindtext}      % 提供 \blindtext 命令，演示用

\title{The Title}
\author{Liam Huang}
\date{\today}
\setcounter{tocdepth}{3}
\begin{document}
\begin{abstract}
This is abstract,this is abstract           % 演示用无意义文字
\begin{keywords}
keyword 1; keyword 2
\end{keywords}

\end{abstract}

\maketitle                  % 打印控制页等
\tableofcontents % 生成目录
\newpage
\section{Introduction}
\label{sec:introduction}
This is an introduction.
This is an introduction.
\subsection{Background}
This is background.



\subsection{Literature Review}
This is the literature review



\subsection{The Description of the Problem}
This is the description of the problem.




\subsection{Our work}
This is our work     
\section{Assumptions}


\textbf{Assumption 1:} The light attenuation and scattering in the underwater environment adhere to the Jaffe-McGlamery imaging model, where light propagations an exponential decay law.

\newpage
\section{Notations}
\setlength{\parindent}{2em} % 设置缩进为2个字符宽度


The core symbols and their definitions used in this study are summarized in Table~\ref{table:notations}, providing an overview of the key parameters and their related meanings.

\begin{table}[ht]
\centering
\caption{Notations used in this literature}
\renewcommand{\arraystretch}{1.5} % Adjust row spacing
\begin{tabular}{>{\centering\arraybackslash}p{0.3\linewidth} >{\centering\arraybackslash}p{0.6\linewidth}} % Two columns with width ratio, centered alignment
\toprule % Top line
\textbf{Symbol} & \textbf{Description} \\ 
\midrule % Middle line

$\text{PSNR}$    & Peak signal-to-noise ratio for image quality evaluation \\

$\text{UCIQE}$   & Underwater Color Image Quality Evaluation metric \\

$\text{UIQM}$    & Underwater Image Quality Measure \\

\bottomrule % Bottom line
\end{tabular}
\label{table:notations}
\end{table}
\newpage
\section{Models}
\subsection{Modeling and Solving of Problem 1}
\subsubsection{Problem Analysis}
\subsubsection{Model Preparation}
\subsubsection{Model Construction}
\subsubsection{Model Solution}
\subsection{Modeling and Solving of Problem 2}
\subsubsection{Problem Analysis}
\subsubsection{Model Preparation}
\subsubsection{Model Construction}
\subsubsection{Model Solution}
\section{Sensitivity Analysis}
\section{Model Evaluation and Promotion}
\subsection{Model Evaluation}
\subsubsection{Advantages}
\subsubsection{Limitations}
\subsection{Future Work}
\subsubsection{Model extension}
\subsubsection{Model application}
\section{Conclusions}

\begin{thebibliography}{9}%宽度9
\bibitem{bib4} L. Prandtl, Fluid motions with very small friction, Proceedings of the 3rd International Mathematical Congress, Heidelberg: H. Schlichting, 1904, 484-491.

\end{thebibliography}
\section*{Appendices}


\end{document}